% from ega1-2-fr.tex
\oldpage[I]{98}

\begin{proposition}[2.1.3]
\label{I.2.1.3-fr}
Si $(X,\mathscr{O}_X)$ est un préschéma, les ensembles ouverts affines forment
une base de la topologie de $X$.
\end{proposition}

En effet, si $V$ est un voisinage ouvert quelconque de $x\in X$, il existe par
hypothèse un voisinage ouvert $W$ de $x$ tel que
$(W,\mathscr{O}_X|W)$ soit un schéma affine~; désignons par $A$ son anneau.
Dans l'espace $W$, $V\cap W$ est un voisinage ouvert de $x$~; donc il existe
$f\in A$ tel que $D(f)$ soit un voisinage ouvert de $x$ contenu dans
$V\cap W$ (\hyperref[I.1.1.10-fr]{1.1.10 (i)}). L'espace annelé
$(D(f),\mathscr{O}_X|D(f))$ est alors un schéma affine d'anneau isomorphe à
$A_f$ (\hyperref[I.1.3.6-fr]{1.3.6}), d'où la proposition.

\begin{proposition}[2.1.4]
\label{I.2.1.4-fr}
L'espace sous-jacent d'un préschéma est un espace de Kolmogoroff.
\end{proposition}

En effet, si $x,y$ sont deux points distincts d'un préschéma $X$, il est
évident qu'il existe un voisinage ouvert de l'un de ces points ne contenant
pas l'autre si $x,y$ ne sont pas dans un même ouvert affine~; et s'ils sont
dans un même ouvert affine, cela résulte de
(\hyperref[I.1.1.8-fr]{1.1.8}).

\begin{proposition}[2.1.5]
\label{I.2.1.5-fr}
Si $(X,\mathscr{O}_X)$ est un préschéma, toute partie fermée irréductible de
$X$ admet un point générique et un seul, et l'application
$x\longmapsto\overline{\{x\}}$ est donc une bijection de $X$ sur l'ensemble
de ses parties fermées irréductibles.
\end{proposition}

En effet, si $Y$ est une partie fermée irréductible de $X$ et $y\in Y$, et si
$U$ est un voisinage ouvert affine de $y$ dans $X$, $U\cap Y$ est partout
dense dans $Y$ et est irréductible
(\hyperref[0.2.1.1-fr]{0, 2.1.1} et
\hyperref[0.2.1.4-fr]{2.1.4})~; donc
(\hyperref[I.1.1.14-fr]{1.1.14}), $U\cap Y$ est l'adhérence dans $U$ d'un
point $x$, et par suite $Y\subset\overline{U}$ est l'adhérence de $x$ dans
$X$. L'unicité du point générique de $X$ résulte de
(\hyperref[I.2.1.4-fr]{2.1.4}) et de
(\hyperref[0.2.1.3-fr]{0, 2.1.3}).

\begin{env}[2.1.6]
\label{I.2.1.6-fr}
Si $Y$ est une partie fermée irréductible de $X$, $y$ son point générique,
l'anneau local $\mathscr{O}_y$ se note aussi $\mathscr{O}_{X/Y}$ et s'appelle
l'\emph{anneau local de $X$ le long de $Y$}, ou l'\emph{anneau local de $Y$
dans $X$}.

Si $X$ lui-même est irréductible et si $x$ est son point générique, on dit
encore que $\mathscr{O}_x$ est l'\emph{anneau des fonctions rationnelles sur
$X$} (cf. \hyperref[section:I.7-fr]{§~7}).
\end{env}

\begin{proposition}[2.1.7]
\label{I.2.1.7-fr}
Si $(X,\mathscr{O}_X)$ est un préschéma, pour toute partie ouverte $U$ de $X$,
l'espace annelé $(U,\mathscr{O}_X|U)$ est un préschéma.
\end{proposition}

Cela résulte aussitôt de la déf. (\hyperref[I.2.1.2-fr]{2.1.2}) et de la prop.
(\hyperref[I.2.1.3-fr]{2.1.3}).

On dit que $(U,\mathscr{O}_X|U)$ est le préschéma \emph{induit} sur $U$ par
$(X,\mathscr{O}_X)$, ou la \emph{restriction} de $(X,\mathscr{O}_X)$ à $U$.

\begin{env}[2.1.8]
\label{I.2.1.8-fr}
On dit qu'un préschéma $(X,\mathscr{O}_X)$ est \emph{irréductible} (resp.
\emph{connexe}) si l'espace sous-jacent $X$ est irréductible (resp. connexe).
On dit qu'un préschéma est \emph{intègre} s'il est \emph{irréductible et
réduit} (cf. (\hyperref[I.5.1.4-fr]{5.1.4})). On dit qu'un préschéma
$(X,\mathscr{O}_X)$ est \emph{localement intègre} si tout $x\in X$ admet un
voisinage ouvert $U$ tel que le préschéma induit sur $U$ par
$(X,\mathscr{O}_X)$ soit intègre.
\end{env}

\subsection{Morphismes de préschémas.}
\label{subsection:I.2.2-fr}
\label{I.2.2-fr}

\begin{definition}[2.2.1]
\label{I.2.2.1-fr}
Étant donnés deux préschémas $(X,\mathscr{O}_X)$, $(Y,\mathscr{O}_Y)$, on
appelle \emph{morphisme (de préschémas)} de $(X,\mathscr{O}_X)$ dans
$(Y,\mathscr{O}_Y)$ tout morphisme d'espaces annelés $(\psi,\theta)$ tel que,
pour tout $x\in X$, $\theta_x^\sharp$ soit un homomorphisme local~:
$\mathscr{O}_{\psi(x)}\longrightarrow\mathscr{O}_x$.
\end{definition}

